% Generated draft: requires mathematical and source review.
\documentclass[11pt]{article}
\usepackage[margin=25mm]{geometry}
\usepackage{fontspec}
\setmainfont{Latin Modern Roman}
\usepackage{amsmath,amssymb,mathtools,unicode-math}
\setmathfont{Latin Modern Math}
\usepackage{enumitem,needspace,xurl,hyperref}
\setlength{\parindent}{0pt}
\setlength{\parskip}{5pt}
\setlength{\emergencystretch}{3em}
\newcommand{\N}{\mathbb{N}}
\newcommand{\Z}{\mathbb{Z}}
\newcommand{\Q}{\mathbb{Q}}
\newcommand{\R}{\mathbb{R}}
\newcommand{\C}{\mathbb{C}}
\newcommand{\F}{\mathbb{F}}
\newcommand{\A}{\mathbb{A}}
\newcommand{\PP}{\mathbb{P}}
\newcommand{\E}{\mathbb{E}}
\newcommand{\eps}{\varepsilon}
\newcommand{\dd}{\,\mathrm{d}}
\begin{document}
\section*{UnsolvedMath ID: 6400001}
\textbf{Source problem code:} catalogue entry
\subsection{Definitions and mathematical statement}
For a smooth knot $K\subset S^3$, let $\Delta_K(t)\in\Z[t,t^{-1}]$ be its Alexander polynomial, considered up to multiplication by $\pm t^j$ unless a normalization is fixed. Equality in the question means equality in this sense, or equality of consistent normalized representatives.

The Fintushel--Stern knot-surgery construction requires a fixed smooth four-manifold $X$, an embedded self-intersection-zero torus $T$ with chosen tubular neighborhood $\nu T\cong T^2\times D^2$, and gluing data. In standard notation,
\[
X_K=\bigl(X\setminus\operatorname{int}(\nu T)\bigr)
\mathop{\cup}_{\varphi_K}\bigl(S^1\times(S^3\setminus\operatorname{int}(\nu K))\bigr).
\]
The boundary identification includes matching the normal meridian of $T$ with the preferred longitude of $K$, together with the remaining framing and orientation data. This describes the type of construction, not a choice of missing data for the source.

For a fixed intended surgery setup, the question is whether
\[
\Delta_K=\Delta_{K'}\quad\Longrightarrow\quad X_K\cong X_{K'}\text{ by a diffeomorphism}
\]
for all knots $K,K'$. The source does not specify the base $X$, torus hypotheses, gluing convention, or whether the diffeomorphism must preserve orientation or a marked complement. These distinctions affect the target.

The supplied summary expects Alexander-polynomial equality not to determine diffeomorphism type and claims same-polynomial distinguishing examples via smooth invariants, but gives no example, invariant, or reference. Those statements are not treated as a verified counterexample to the exact fixed-setup question. Equality of one knot-derived invariant cannot by itself be substituted for a classification theorem or its negation.
\subsection{Short English statement}
For a fixed Fintushel--Stern knot-surgery construction, does equality of Alexander polynomials force the resulting four-manifolds to be diffeomorphic? The surgery setup and scope of the supplied distinguishing claims need verification.
\subsection{Sources}
\begin{itemize}
\item Local UnsolvedMath registry, permanent ID 6400001, exact source code catalogue entry, ``Distinguishing Fintushel--Stern Manifolds''; proposed by Simon Donaldson, 2008. Registry status partially solved, classification PARTIAL-PROGRESS.
\item Simon Donaldson, \emph{Some Problems in Differential Geometry and Topology} (2008), explicit question/program 1, \url{https://wwwf.imperial.ac.uk/~skdona/NONLINEARITYPROBLEMS.PDF}. Supplied normalized pdf-text extraction accessed July 29, 2026. No PDF was accessed here.
\item The supplied research summary describes expectations and claims non-diffeomorphic, often homeomorphic, knot-surgery examples for knots sharing an Alexander polynomial. It supplies no precise construction, base manifold, invariant, or publication supporting that claim, so its applicability to the exact source question remains unresolved.
\item Provenance says the author survey presents the programme as outstanding and flags current status and statement-reuse terms NEEDS\_REVIEW; the document is publicly author-hosted.
\end{itemize}
\end{document}
